\chapter{Linux \& Customizing my Computer}
Now that I am dual booting my computer I am going to make plans for customizing it and making sure I have all of the things I want installed.i

\section{Introductory Plans}
I want the libraries, package management tools, set up emacs, Terminal Multiplexing, compililers, toolchain, git, setting up Windows/Linux workflow,
and automation
\section{Package Manegment}
ripgreg, fd, bat, fzf

\section{Terminal Emulator}

I am thinking of choicing alacritty

\section{Editor}
While I plan to put all of my Emacs stuff in the GitHub to transer, what I am thinking
is for this time to use NeoVim and then from there decide which to use.

\section{AI summary}

Rather than write the rest of it I had ChatGPT write a summary:

\subsection{Core System Setup}

\subsubsection{Installation}

\begin{lstlisting}
sudo apt update && sudo apt upgrade -y

sudo apt install -y \
build-essential cmake make ninja-build \
git curl wget \
htop tmux openssh-client openssh-server
\end{lstlisting}

\subsubsection{Purpose}

\begin{itemize}
    \item \texttt{build-essential}: Provides compilers and standard libraries
    \item \texttt{cmake}: Cross-platform build system
    \item \texttt{tmux}: Persistent sessions for long simulations
    \item \texttt{ssh}: Remote access for distributed workflows
\end{itemize}

\subsection{Compilers}

\subsubsection{Installation}

\begin{lstlisting}
sudo apt install -y gcc g++ gfortran clang
\end{lstlisting}

\subsubsection{Why Multiple Compilers?}

Different compilers provide:

\begin{itemize}
    \item Different optimization strategies
    \item Better compatibility with legacy codes
    \item Access to vectorization and architecture-specific tuning
\end{itemize}

\subsubsection{Usage Example}

\begin{lstlisting}
gfortran -O3 -march=native simulation.f90 -o sim
\end{lstlisting}

\subsection{Scientific Libraries}

\subsubsection{Installation}

\begin{lstlisting}
sudo apt install -y \
libopenblas-dev liblapack-dev \
libfftw3-dev \
libhdf5-dev \
libopenmpi-dev openmpi-bin
\end{lstlisting}

\subsubsection{Mathematical Role}

\begin{itemize}
    \item \textbf{BLAS/LAPACK}: Linear algebra (matrix solves, eigenvalues)
    \item \textbf{FFTW}: Fast Fourier transforms
    \item \textbf{HDF5}: Structured scientific data storage
    \item \textbf{MPI}: Parallel distributed computation
\end{itemize}

\subsubsection{MPI Example}

\begin{lstlisting}
mpirun -np 4 ./simulation
\end{lstlisting}

\subsection{Python Scientific Environment}

\subsubsection{Installation (Conda)}

\begin{lstlisting}
conda create -n physics python=3.11
conda activate physics

conda install numpy scipy matplotlib pandas numba sympy h5py
\end{lstlisting}

\subsubsection{Role in Workflow}

Python acts as:

\begin{itemize}
    \item Data analysis layer
    \item Visualization interface
    \item Rapid prototyping environment
\end{itemize}

\subsubsection{Example}

\begin{lstlisting}
import h5py
import numpy as np

f = h5py.File("output.h5", "r")
data = np.array(f["density"])
\end{lstlisting}

\subsection{Fortran Development}

\subsubsection{Tools}

\begin{lstlisting}
pip install fpm
\end{lstlisting}

\subsubsection{Why Fortran?}

\begin{itemize}
    \item Efficient array operations
    \item Strong numerical performance
    \item Widely used in legacy HPC codes
\end{itemize}

\subsection{Visualization}

\subsubsection{Installation}

\begin{lstlisting}
sudo apt install gnuplot paraview
\end{lstlisting}

\subsubsection{Purpose}

\begin{itemize}
    \item \textbf{gnuplot}: Quick plotting
    \item \textbf{ParaView}: Large-scale simulation visualization
\end{itemize}

\subsection{LaTeX Environment}

\subsubsection{Installation}

\begin{lstlisting}
sudo apt install texlive-full
\end{lstlisting}

\subsubsection{Purpose}

LaTeX is used for:

\begin{itemize}
    \item Research papers
    \item Technical documentation
    \item Mathematical derivations
\end{itemize}

\subsection{Editors}

\subsubsection{Installation}

\begin{lstlisting}
sudo apt install emacs
\end{lstlisting}

\subsubsection{Usage Philosophy}

\begin{itemize}
    \item Emacs: Integrated research environment
\end{itemize}

\subsection{Containers and Reproducibility}

\subsubsection{Installation}

\begin{lstlisting}
sudo apt install docker.io
\end{lstlisting}

\subsubsection{Purpose}

\begin{itemize}
    \item Reproducible environments
    \item Dependency isolation
    \item Portable experiments
\end{itemize}

\subsection{Workflow Architecture}

\subsubsection{Full Pipeline}

\begin{enumerate}
    \item Write simulation code (Fortran/C++)
    \item Compile with optimized libraries
    \item Run locally or via MPI
    \item Output data in HDF5 format
    \item Analyze using Python
    \item Visualize results
    \item Document in LaTeX
\end{enumerate}

\subsection{Directory Structure}

\begin{lstlisting}
~/projects/
~/simulations/
~/data/
~/papers/
~/notes/
\end{lstlisting}

\subsection{Advanced Topics}

\subsection{Performance Optimization}

\begin{itemize}
    \item Use \texttt{-O3 -march=native}
    \item Profile with \texttt{perf} or \texttt{gprof}
    \item Optimize memory access patterns
\end{itemize}

\subsubsection{Parallel Scaling}

\begin{itemize}
    \item Domain decomposition
    \item MPI communication patterns
    \item Load balancing
\end{itemize}

\subsubsection{Reproducibility Strategy}

\begin{itemize}
    \item Version control (git)
    \item Environment locking (conda)
    \item Containerization (Docker)
\end{itemize}


\section{3/20/26}

I had some trouble with having different scalling for different displays. I ended up figuring out to change the GNOME from X11 to Wayland.

I was not able to make sure that the cursor switching works as I wanted, as far as I can tell there is nothing I can do to fix that.
